\section{Tessuto muscolare}

\textbf{Tessuto muscolare}: cellule con la capacità di contrarsi, mediante miofilamenti di actina e miosina. Il meccanismo di contrazione è simile in tutti e tre i tipi di tessuto muscolare, ma le cellule muscolari sono diverse nella loro organizzazione interna.

\begin{itemize}
  \item \textbf{striato} (o scheletrico): tutti i muscoli che muovono o stabilizzano lo scheletro; muscoli che controllano le vie di ingresso e di uscita dei tratti digerente, respiratorio e urinario;
  \item \textbf{cardiaco} (o miocardio): cuore;
  \item \textbf{liscio} (o viscerale): pareti dei vasi sanguigni, organi degli apparati respiratorio e urogenitale.
\end{itemize}

\rubrica{Caratteristiche}
\begin{itemize}
  \item \textbf{eccitabilità}: risponde alla stimolazione di nervi e ormoni $\rightarrow$ contrazione;
  \item \textbf{elasticità}: se stirato, ritorna alla sua lunghezza di riposo;
  \item \textbf{contrattilità ed estensibilità}: può accorciarsi o estendersi rispetto alla normale lunghezza di riposo, con produzione di forza.
\end{itemize}

% ---------------------------------------------------------------------------
\subsection{Tessuto muscolare scheletrico (striato volontario)}

\begin{itemize}
  \item cellule molto grandi, cilindriche: \textbf{fibrocellule} o \textbf{fibre muscolari};
  \item incapaci di dividersi e di proliferare;
  \item cellule unite a formare un \textbf{sincizio};
  \item costituito da 100 fino a 100.000 cellule multinucleate allungate, con nuclei in posizione periferica;
  \item l'unità contrattile è il \textbf{sarcomero}.
\end{itemize}

\rubrica{Involucri connettivali}
\begin{itemize}
  \item \textbf{epimisio}: avvolge il muscolo (organo);
  \item \textbf{perimisio}: avvolge il fascicolo muscolare (fascio di fibre);
  \item \textbf{endomisio}: avvolge la singola fibra muscolare (cellula).
\end{itemize}
La singola fibra comprende sarcolemma, sarcoplasma, miofibrille, mitocondri, cellula satellite, nucleo, ed è servita da capillari.

\rubrica{Cellule satelliti}
\begin{itemize}
  \item \textbf{mioblasti quiescenti};
  \item conferiscono capacità rigenerativa al tessuto;
  \item localizzate tra il sarcolemma di una fibra muscolare e la lamina basale circostante.
\end{itemize}

\rubrica{Microcircolo}
\begin{itemize}
  \item \textbf{arteriola}: trasporta ossigeno e nutrienti;
  \item \textbf{capillari}: le pareti sottili consentono la diffusione; 3--5 per fibra;
  \item \textbf{venula}: trasporta il sangue refluo.
\end{itemize}

\rubrica{Sistema di membrane}
\begin{itemize}
  \item \textbf{sarcolemma}: membrana plasmatica della fibra;
  \item \textbf{tubulo trasverso T}: invaginazione del sarcolemma;
  \item \textbf{reticolo sarcoplasmatico}: cisterne fenestrate e \textbf{cisterne terminali};
  \item \textbf{triade}: un tubulo T con le due cisterne terminali affiancate;
  \item \textbf{filamenti intermedi} (desmina, vimentina, sinemina), disco Z, mitocondri.
\end{itemize}

% ---------------------------------------------------------------------------
\subsection{Il sarcomero}

Unità contrattile della miofibrilla, delimitata da due \textbf{linee Z} (o dischi Z). Comprende:
\begin{itemize}
  \item \textbf{banda A}, \textbf{banda H}, \textbf{banda I};
  \item \textbf{linea M}: al centro del sarcomero;
  \item \textbf{filamenti sottili} (actina) e \textbf{filamenti spessi} (miosina), collegati dai \textbf{ponti trasversali};
  \item proteine strutturali (titina, nebulina) e di ancoraggio.
\end{itemize}

\rubrica{Filamenti sottili}
\begin{itemize}
  \item \textbf{actina}: proteina globulare (actina G) che polimerizza a formare le catene di actina F (due catene di 300--400 molecole globulari); ciascun monomero di actina G presenta un \textbf{sito attivo} per la miosina, in condizioni di riposo mascherato dalla tropomiosina;
  \item \textbf{tropomiosina}: proteina filamentosa che copre i siti attivi dell'actina, impedendo l'interazione con la miosina;
  \item \textbf{troponina}: proteina globulare che si lega alle due precedenti per mantenerne i rapporti:
  \begin{itemize}
    \item \textbf{TnC}: lega il $Ca^{2+}$;
    \item \textbf{TnI}: subunità inibitrice;
    \item \textbf{TnT}: subunità che lega la tropomiosina.
  \end{itemize}
\end{itemize}
Alle estremità del filamento sottile: \textbf{tropomodulina} (estremità libera) e \textbf{Cap Z}, in corrispondenza della linea Z.

\rubrica{Filamenti spessi}
\begin{itemize}
  \item \textbf{miosina}: proteina filamentosa formata da una porzione allungata (\textbf{coda}) e da un'estremità globosa (\textbf{testa}); tra coda e testa un \textbf{punto di flessione} (collo);
  \item la testa presenta il sito di legame dell'actina e il sito di legame dell'ATP (interagisce con actina e ATP);
  \item struttura: doppia elica di \textbf{catene pesanti} (coda) e \textbf{catene leggere} (testa);
  \item circa 250 molecole di miosina per filamento spesso, con una \textbf{area nuda} centrale.
\end{itemize}

\rubrica{Disposizione delle miosine}
Nel filamento spesso le molecole di miosina sono disposte con \textbf{polarità opposta}, essenziale per il meccanismo della contrazione. La parte centrale del filamento è priva di teste (ponti trasversi) per un'ampiezza di 2500\,\AA.

\rubrica{Proteine strutturali e di ancoraggio}
\begin{itemize}
  \item \textbf{titina}: dominio elastico a livello della banda I, dominio inestensibile a livello della banda A;
  \item \textbf{nebulina}: associata ai filamenti sottili;
  \item \textbf{$\alpha$-actinina}: alla linea Z;
  \item \textbf{proteine di ancoraggio}: \textbf{distrofina} (connette la F-actina al sarcolemma), \textbf{desmina} (connette la miofibrilla, a livello della linea Z, al sarcolemma).
\end{itemize}

% ---------------------------------------------------------------------------
\subsection{Tessuto muscolare cardiaco (miocardio, striato involontario)}

\begin{itemize}
  \item cellule corte e ramificate, di dimensioni ridotte;
  \item incapaci di dividersi e di proliferare;
  \item mononucleate;
  \item presenza di filamenti di actina e di miosina;
  \item ampi collegamenti le une con le altre, mediante \textbf{dischi intercalari};
  \item membrane saldate insieme da desmosomi e \textit{gap junctions};
  \item presenza di \textbf{cellule pacemaker} (segnapasso);
  \item situato nel cuore, in rapporto con tessuto connettivo (valvole cardiache).
\end{itemize}

\rubrica{Funzioni}
\begin{itemize}
  \item circolazione sanguigna;
  \item mantenimento della pressione del sangue (idrostatica).
\end{itemize}

\rubrica{Disco intercalare}
Zona di giunzione tra cardiociti adiacenti, comprende:
\begin{itemize}
  \item \textbf{giunzione comunicante} (\textit{gap junction});
  \item \textbf{fascia aderente};
  \item \textbf{desmosomi}.
\end{itemize}

% ---------------------------------------------------------------------------
\subsection{Tessuto muscolare liscio (viscerale involontario)}

\begin{itemize}
  \item cellule piccole e fusiformi;
  \item capacità di suddividersi e di proliferare;
  \item mononucleate;
  \item presenza di filamenti di actina e di miosina, ma assenza della caratteristica striatura;
  \item si contrae indipendentemente dalla volontà, grazie al sistema nervoso autonomo.
\end{itemize}
Presente nei muscoli viscerali (tonaca muscolare) degli organi cavi: tubo digerente, albero respiratorio, vie urinarie e genitali, arterie e vene.

\rubrica{Funzioni}
\begin{itemize}
  \item progressione di cibo, feci, urina e secrezioni;
  \item controllo del calibro delle vie respiratorie;
  \item controllo del calibro dei vasi sanguigni.
\end{itemize}

\rubrica{Organizzazione interna}
\begin{itemize}
  \item \textbf{corpi densi} collegati da filamenti intermedi di desmina;
  \item sarcomeri assenti;
  \item le cisterne del reticolo sarcoplasmatico sono rimpiazzate dalle \textbf{caveole};
  \item la cellula contratta assume una forma rotondeggiante e bernoccoluta.
\end{itemize}

\rubrica{Innervazione}
\begin{itemize}
  \item \textbf{muscolo liscio unitario}: innervazione unitaria;
  \item \textbf{muscolo liscio multiunitario}: innervazione multiunitaria.
\end{itemize}

\rubrica{Meccanismo della contrazione}
\begin{enumerate}
  \item gli ioni $Ca^{2+}$ entrano nel citosol dal fluido extracellulare, attraverso canali per il $Ca^{2+}$ voltaggio-dipendenti o -indipendenti, oppure dallo scarso reticolo sarcoplasmatico;
  \item il $Ca^{2+}$ si lega alla \textbf{calmodulina} e la attiva;
  \item la calmodulina attivata attiva a sua volta la \textbf{chinasi della catena leggera della miosina};
  \item la chinasi attivata catalizza il trasferimento del fosfato sulla miosina, attivando l'ATPasi miosinica (stato inattivo = catene leggere non fosforilate; stato attivo = catene leggere fosforilate);
  \item la miosina attivata forma ponti trasversali con l'actina dei miofilamenti sottili, dando inizio all'accorciamento.
\end{enumerate}

% ---------------------------------------------------------------------------
\subsection{Confronto tra i tre tipi di tessuto muscolare}

\begin{table}[htbp]
  \centering
  \small
  \renewcommand{\arraystretch}{1.3}
  \begin{tabularx}{\textwidth}{|l|X|X|X|}
    \hline
    \textbf{Caratteristica} & \textbf{Scheletrico} & \textbf{Cardiaco} & \textbf{Liscio} \\
    \hline
    Componenti connettivali & epimisio, perimisio ed endomisio & endomisio, attaccato allo scheletro fibroso del cuore & endomisio \\
    \hline
    Miofibrille costituite da sarcomeri & sì & sì, ma le miofibrille hanno spessore irregolare & no, ma i filamenti di actina e miosina sono presenti \\
    \hline
    Tubuli T e sito dell'invaginazione & sì, due per sarcomero a livello delle giunzioni A--I & sì, a livello delle linee Z; diametro maggiore rispetto ai tubuli del muscolo scheletrico & no, caveole lungo il sarcolemma \\
    \hline
    Reticolo sarcoplasmatico elaborato & sì & meno del muscolo scheletrico; poche cisterne terminali & come nel cardiaco; una parte del RS è in contatto con il sarcolemma \\
    \hline
    Giunzioni comunicanti & no & sì, a livello dei dischi intercalari & sì, nel muscolo unitario \\
    \hline
    Giunzioni neuromuscolari individuali & sì & no & no nel muscolo unitario; sì nel multiunitario \\
    \hline
    Fonte di $Ca^{2+}$ per la contrazione & reticolo sarcoplasmatico & RS e fluido extracellulare & RS e fluido extracellulare \\
    \hline
  \end{tabularx}
\end{table}
